(1) 高折射率侧的波动方程给 \(k_z=\sqrt{n_1^2k_0^2-k_x^2}=n_1k_0\cos\theta>0\)。低折射率侧因 \(k_x>n_2k_0\)，纵向波数为 \(i\kappa\)，\(\kappa=\sqrt{k_x^2-n_2^2k_0^2}>0\)；透入低折射率介质的单个倏逝分量随 \(z\) 按 \(e^{-\kappa z}\) 下降。
(2) TE电场切向分量与其法向导数在非磁性界面连续，故 \(1+r_s=t_s\)、\(k_z(1-r_s)=i\kappa t_s\)，解得 \(r_s=(k_z-i\kappa)/(k_z+i\kappa)\)。分子分母互为复共轭，因此 \(|r_s|=1\)；在连续相位分支下 \(\phi_s=-2\arctan(\kappa/k_z)\)，虽然反射功率为一，相位仍随角度改变。
(3) 反射束的角谱积分含相位 \(k_xx+\phi_s(k_x)\)，在中心波数驻相处有 \(x+\partial\phi_s/\partial k_x=0\)，所以 \(\Delta_{\rm GH}=-\partial\phi_s/\partial k_x\)。利用 \(dk_z/dk_x=-k_x/k_z\)、\(d\kappa/dk_x=k_x/\kappa\)，得 \(d\phi_s/dk_x=-2k_x/(k_z\kappa)\)，从而 \(\Delta_{\rm GH}=2k_x/(k_z\kappa)=2\tan\theta/\kappa>0\)。
(4) 当 \(\theta\downarrow\theta_c=\arcsin(n_2/n_1)\) 时 \(\kappa\to0^+\)，简单驻相预测位移增强；实际有限宽束包含临界角两侧谱分量，不能将局部导数发散外推为无限观测位移。\(\theta\to\pi/2\) 时 \(k_z\to0\)，相位对角谱也非常敏感，有限束孔径、偏振、衍射及非平面束条件限制单色窄谱近似的有效性。
